% year: תשפ"ב
% type: מבחן
% semester: א
% moed: ב
% number: 2א
% points: 8
% categories: func-analysis
% source: exams/מבחנים/תשפב/מועד ב תשפב.pdf

\begin{question}
(8 נק') מצאו נקודות קיצון מקומי ותחומי עלייה וירידה של הפונקציה
\[ f(x)=\frac{2x}{x^2+1} \]
\end{question}

\begin{hint}
חשבו את $f'$ לפי כלל המנה ובדקו את סימנה; המכנה של $f'$ חיובי תמיד.
\end{hint}

\begin{solution}
$f$ מוגדרת וגזירה בכל $\R$ (המכנה חיובי). לפי כלל המנה:
\[ f'(x)=\frac{2(x^2+1)-2x\cdot2x}{(x^2+1)^2}=\frac{2(1-x^2)}{(x^2+1)^2}. \]
$f'(x)=0\iff x=\pm1$. סימן $f'$ הוא סימן $1-x^2$:
\begin{itemize}
  \item $f'<0$ ב-$(-\infty,-1)$: $f$ \textbf{יורדת}.
  \item $f'>0$ ב-$(-1,1)$: $f$ \textbf{עולה}.
  \item $f'<0$ ב-$(1,\infty)$: $f$ \textbf{יורדת}.
\end{itemize}
לכן ב-$x=-1$ יש \textbf{מינימום מקומי} $f(-1)=-1$, וב-$x=1$ יש \textbf{מקסימום מקומי} $f(1)=1$.
(למעשה אלו גם קיצון מוחלט: $|2x|\le x^2+1$ כי $(|x|-1)^2\ge0$, ולכן $-1\le f\le 1$.)
\end{solution}
